\documentclass{report}
\usepackage{luatexja} % LuaTeXで日本語を使うためのパッケージ
\usepackage{luatexja-fontspec} % LuaTeX用の日本語フォント設定

% --- 数学関連 ---
\usepackage{amsmath, amssymb, amsfonts, mathtools, bm, amsthm} % 基本的な数学パッケージ
\usepackage{type1cm, upgreek} % 数式フォントとギリシャ文字
\usepackage{physics, mhchem} % 物理や化学の記号や式の表記を簡単にする

% --- 表関連 ---
\usepackage{multirow, longtable, tabularx, array, colortbl, dcolumn, diagbox} % 表のレイアウトを柔軟にする
\usepackage{tablefootnote, truthtable} % 表中に注釈を追加、真理値表
\usepackage{tabularray} % 高度な表組みレイアウト

% --- グラフィック関連 ---
\usepackage{tikz, graphicx} % 図の描画と画像の挿入
\usepackage{background} % ウォーターマークの設定
\usepackage{caption, subcaption} % 図や表のキャプション設定
\usepackage{float, here} % 図や表の位置指定

% --- レイアウトとページ設定 ---
\usepackage{fancyhdr} % ページヘッダー、フッター、余白の設定
\usepackage[top = 20truemm, bottom = 20truemm, left = 20truemm, right = 20truemm]{geometry}
\usepackage{fancybox, ascmac} % ボックスのデザイン

% --- 色とスタイル ---
\usepackage{xcolor, color, colortbl, tcolorbox} % 色とカラーボックス
\usepackage{listings, jvlisting} % コードの色付けとフォーマット

% --- 参考文献関連 ---
\usepackage{biblatex, usebib} % 参考文献の管理と挿入
\usepackage{url, hyperref} % URLとリンクの設定

% --- その他の便利なパッケージ ---
\usepackage{footmisc} % 脚注のカスタマイズ
\usepackage{multicol} % 複数段組
\usepackage{comment} % コメントアウトの拡張
\usepackage{siunitx} % 単位の表記
\usepackage{docmute}
% \usepackage{appendix}
% --- tcolorboxとtikzの設定 ---
\tcbuselibrary{theorems, breakable} % 定理のボックスと改ページ設定
\usetikzlibrary{decorations.markings, arrows.meta, calc} % tikzの装飾や矢印の設定

% --- 定理スタイルと数式設定 ---
\theoremstyle{definition} % 定義スタイル
\numberwithin{equation}{section} % 式番号をサブセクション単位でリセット

% --- hyperrefの設定 ---
\hypersetup{
  setpagesize = false,
  bookmarks = true,
  bookmarksdepth = tocdepth,
  bookmarksnumbered = true,
  colorlinks = false,
  pdftitle = {}, % PDFタイトル
  pdfsubject = {}, % PDFサブジェクト
  pdfauthor = {}, % PDF作者
  pdfkeywords = {} % PDFキーワード
}

% --- siunitxの設定 ---
\sisetup{
  table-format = 1.5, % 小数点以下の桁数
  table-number-alignment = center, % 数値の中央揃え
}

% --- 透かし画像の設定 --- 
\backgroundsetup{
  scale=0.5,                       % 画像のスケール
  % color=black,                   % 画像の色（透かし用に半透明が推奨）
  opacity=0.2,                   % 透かしの透明度（0が完全透明、1が完全不透明）
  angle=0,                       % 画像の角度
  position = current page.south east,  % ページの右下
  hshift=-6cm, % 右方向へのシフト（負の値で内側に移動）
  vshift=5cm,  % 上方向へのシフト（正の値で内側に移動）
  contents={\includegraphics{./fig/appilogo-circular-full.png}} % 画像のパス
}

% --- その他の設定 ---
\allowdisplaybreaks % 数式の途中改ページ許可
\newcolumntype{t}{!{\vrule width 0.1pt}} % 新しいカラムタイプ
\newcolumntype{b}{!{\vrule width 1.5pt}} % 太いカラム
\UseTblrLibrary{amsmath, booktabs, counter, diagbox, functional, hook, html, nameref, siunitx, varwidth, zref} % tabularrayのライブラリ
\setlength{\columnseprule}{0.4pt} % カラム区切り線の太さ
\captionsetup[figure]{font = bf} % 図のキャプションの太字設定
\captionsetup[table]{font = bf} % 表のキャプションの太字設定
\captionsetup[lstlisting]{font = bf} % コードのキャプションの太字設定
\captionsetup[subfigure]{font = bf, labelformat = simple} % サブ図のキャプション設定
\setcounter{secnumdepth}{4} % セクションの深さ設定
\newcolumntype{d}{D{.}{.}{5}} % 数値のカラム
\newcolumntype{M}[1]{>{\centering\arraybackslash}m{#1}} % センター揃えのカラム
\DeclareMathOperator{\diag}{diag}
\everymath{\displaystyle} % 数式のスタイル
\newcommand{\inner}[2]{\left\langle #1, #2 \right\rangle}
\renewcommand{\figurename}{図}
\renewcommand{\i}{\mathrm{i}} % 複素数単位i
\renewcommand{\laplacian}{\grad^2} % ラプラシアンの記号
\renewcommand{\thesubfigure}{(\alph{subfigure})} % サブ図の番号形式
\newcommand{\m}[3]{\multicolumn{#1}{#2}{#3}} % マルチカラムのショートカット
\renewcommand{\r}[1]{\mathrm{#1}} % mathrmのショートカット
\newcommand{\e}{\mathrm{e}} % 自然対数の底e
\newcommand{\Ef}{E_{\mathrm{F}}} % フェルミエネルギー
\renewcommand{\c}{\si{\degreeCelsius}} % 摂氏記号
\renewcommand{\d}{\r{d}} % d記号
\renewcommand{\t}[1]{\texttt{#1}} % タイプライタフォント
\newcommand{\kb}{k_{\mathrm{B}}} % ボルツマン定数
% \renewcommand{\phi}{\varphi} % ϕをφに変更
\renewcommand{\epsilon}{\varepsilon}
\newcommand{\fullref}[1]{\textbf{\ref{#1} \nameref{#1}}}
\newcommand{\reff}[1]{\textbf{図\ref{#1}}} % 図参照のショートカット
\newcommand{\reft}[1]{\textbf{表\ref{#1}}} % 表参照のショートカット
\newcommand{\refe}[1]{\textbf{式\eqref{#1}}} % 式参照のショートカット
\newcommand{\refp}[1]{\textbf{コード\ref{#1}}} % コード参照のショートカット
\renewcommand{\lstlistingname}{コード} % コードリストの名前
\renewcommand{\theequation}{\thesection.\arabic{equation}} % 式番号の形式
\renewcommand{\footrulewidth}{0.4pt} % フッターの線
\newcommand{\mar}[1]{\textcircled{\scriptsize #1}} % 丸囲み文字
\newcommand{\combination}[2]{{}_{#1} \mathrm{C}_{#2}} % 組み合わせ
\newcommand{\thline}{\noalign{\hrule height 0.1pt}} % 細い横線
\newcommand{\bhline}{\noalign{\hrule height 1.5pt}} % 太い横線

% --- カスタム色定義 ---
\definecolor{burgundy}{rgb}{0.5, 0.0, 0.13} % バーガンディ色
\definecolor{charcoal}{rgb}{0.21, 0.27, 0.31} % チャコール色
\definecolor{forest}{rgb}{0.0, 0.35, 0} % 森の緑色

% --- カスタム定理環境の定義 ---
\newtcbtheorem[number within = chapter]{myexc}{練習問題}{
  fonttitle = \gtfamily\sffamily\bfseries\upshape,
  colframe = forest,
  colback = forest!2!white,
  rightrule = 1pt,
  leftrule = 1pt,
  bottomrule = 2pt,
  colbacktitle = forest,
  theorem style = standard,
  breakable,
  arc = 0pt,
}{exc-ref}
\newtcbtheorem[number within = chapter]{myprop}{命題}{
  fonttitle = \gtfamily\sffamily\bfseries\upshape,
  colframe = blue!50!black,
  colback = blue!50!black!2!white,
  rightrule = 1pt,
  leftrule = 1pt,
  bottomrule = 2pt,
  colbacktitle = blue!50!black,
  theorem style = standard,
  breakable,
  arc = 0pt
}{proposition-ref}
\newtcbtheorem[number within = chapter]{myrem}{注意}{
  fonttitle = \gtfamily\sffamily\bfseries\upshape,
  colframe = yellow!20!black,
  colback = yellow!50,
  rightrule = 1pt,
  leftrule = 1pt,
  bottomrule = 2pt,
  colbacktitle = yellow!20!black,
  theorem style = standard,
  breakable,
  arc = 0pt
}{remark-ref}
\newtcbtheorem[number within = chapter]{myex}{例題}{
  fonttitle = \gtfamily\sffamily\bfseries\upshape,
  colframe = black,
  colback = white,
  rightrule = 1pt,
  leftrule = 1pt,
  bottomrule = 2pt,
  colbacktitle = black,
  theorem style = standard,
  breakable,
  arc = 0pt
}{example-ref}
\newtcbtheorem[number within = chapter]{exc}{Requirement}{myexc}{exc-ref}
\newcommand{\rqref}[1]{{\bfseries\sffamily 練習問題 \ref{exc-ref:#1}}}
\newtcbtheorem[number within = chapter]{definition}{Definition}{mydef}{definition-ref}
\newcommand{\dfref}[1]{{\bfseries\sffamily 定義 \ref{definition-ref:#1}}}
\newtcbtheorem[number within = chapter]{prop}{命題}{myprop}{proposition-ref}
\newcommand{\prref}[1]{{\bfseries\sffamily 命題 \ref{proposition-ref:#1}}}
\newtcbtheorem[number within = chapter]{rem}{注意}{myrem}{remark-ref}
\newcommand{\rmref}[1]{{\bfseries\sffamily 注意 \ref{remark-ref:#1}}}
\newtcbtheorem[number within = chapter]{ex}{例題}{myex}{example-ref}
\newcommand{\exref}[1]{{\bfseries\sffamily 例題 \ref{example-ref:#1}}}
% --- 再定義コマンド ---
% \mathtoolsset{showonlyrefs=true} % 必要な式番号のみ表示
\pagestyle{fancy} % ヘッダー・フッターのスタイル設定
\chead{応用量子物性講義ノート} % 中央ヘッダー
% \rhead{}
\fancyhead[R]{\rightmark}
\renewcommand{\sectionmark}[1]{\markright{\thesection\ #1}}
\cfoot{\thepage} % 中央フッターにページ番号
\lhead{}
\rfoot{Yuto Masuda and Haruki Aoki} % 右フッターに名前
\setcounter{tocdepth}{4} % 目次の深さ
\makeatletter
\@addtoreset{equation}{section} % サブセクションごとに式番号をリセット
\makeatother

% --- メタ情報 ---
\title{応用量子物性講義ノート}
\date{更新日\today}
\author{Yuto Masuda and Haruki Aoki}

\begin{document}
  散乱体が球対称ポテンシャル$V(r)$を持つとする．Schrödinger方程式は以下のようになる．
  \begin{align}
    \label{SE_Scattering}
    \qty(-\frac{\hbar^2}{2m}\nabla^2 + V(r))\psi(\bm{r}) = E \psi(\bm{r})
  \end{align}
  ここで，$V(r)$は$r \to \infty$で十分早く$V \to 0$となるとする．$z$軸に沿う入射波は平面波なので
  \begin{align}
    \psi_{\r{in}} = \e^{\i k z}\ (z \to \infty)
  \end{align}
  と表せる．また，散乱波は外向きの球面波となるので
  \begin{align}
    \psi_{\r{sc}} \simeq \frac{\e^{\i kr}}{r}\ (r \to \infty)
  \end{align}
  である．散乱問題とは，$\r \to \infty$で
  \begin{align}
    \psi(\bm{r}) = \e^{\i kz} + f(\theta) \frac{\e^{\i kr}}{r}
  \end{align}
  を満たす式(\ref{SE_Scattering})の定常解を求めることである．$f(\theta)$を散乱振幅という．

  この問題を解くための準備として確率密度($\rho = \psi^{*} \psi$)の時間変化を考える．
  \begin{align}
    \frac{\partial}{\partial t}\rho &= \frac{\partial}{\partial t}(\psi^{*}\psi)\\
    &= \frac{\partial \psi^{*}}{\partial t}\psi + \psi^{*}\frac{\partial \psi}{\partial t}\\
    &= -\frac{1}{\i\hbar}\qty(-\frac{\hbar^2}{2m}\nabla^2 + V)\psi^{*}\psi + \psi^{*}\frac{1}{\i\hbar}\qty(-\frac{\hbar^2}{2m}\nabla^2 + V)\psi\\
    &= \frac{\hbar}{2m\i}\qty[(\nabla^2 \psi^{*})\psi - \psi^{*}(\nabla^2\psi)]\\
    &= -\frac{\hbar}{2m\i}\nabla \cdot (\psi^{*}\nabla \psi - \psi \nabla \psi^{*})
  \end{align}
  ここで，確率流密度を
  \begin{align}
    \bm{j} &= \frac{\hbar}{2m\i} (\psi^{*}\nabla \psi - \psi \nabla \psi^{*})\\
    &= \frac{\hbar}{m}\Im (\psi^{*}\nabla\psi)
  \end{align}
  と定義すれば，確率密度に対する連続の式
  \begin{align}
    \frac{\partial}{\partial t} \rho = -\nabla \cdot \bm{j}
  \end{align}
  を得る．

  次に微分断面積と散乱振幅の関係を考える．
  入射波は$\psi_{\r{in}} = \e^{\i kz}$である．入射波の確率流密度は
  \begin{align}
    j_z &= \frac{\hbar}{m} \Im (\e^{-\i kz} \i k e^{\i kz}) \\
    &= \frac{\hbar k}{m}
  \end{align}
  である．散乱波は$\psi_{\r{sc}} = \frac{f(\theta)}{r}\e^{\i kr}$である．散乱波の確率流密度は
  \begin{align}
    j_r(\theta) &= \frac{\hbar}{m}\Im (\psi^{*}_{\r{sc}}\nabla\psi_{\r{sc}})\\
    &\simeq \frac{\hbar}{m}\Im (\psi^{*}_{\r{sc}}\frac{\partial}{\partial r}\psi_{\r{sc}})\\
    &= \frac{\hbar}{m} \Im \qty[\frac{f(\theta)}{r} \e^{-\i kr}\qty(\frac{\i k}{r} - \frac{1}{r^2})f(\theta)\e^{\i kr}]\\
    &\simeq \frac{\hbar}{m} \Im \qty[\frac{f(\theta)}{r} \e^{-\i kr}\frac{\i k}{r}f(\theta)\e^{\i kr}]\\
    &= \frac{\hbar}{m}\frac{k}{r^2} \abs{f(\theta)}^2\\
    &= \frac{\abs{f(\theta)}^2}{r^2}j_z
  \end{align}
  である．ここで，1行目から2行目では
  \begin{align}
    \nabla \psi_{\r{sc}} &= \qty(\frac{\partial}{\partial r}\bm{e}_r + \frac{1}{r}\frac{\partial}{\partial \theta}\bm{e}_{\theta} + \frac{1}{r \sin \theta}\frac{\partial}{\partial \phi} \bm{e}_{\phi})\psi_{\r{sc}}\\
    &\simeq \frac{\partial}{\partial r}\psi_{\r{sc}} \bm{e}_r\ (r \to \infty)
  \end{align}
  という近似を用いた．
  
  微分断面積と粒子数の関係
  \begin{align}
    \dd{N} = \sigma (\theta) N \dd{\Omega}
  \end{align}
  の両辺を$N$で割る．
  \begin{align}
    \frac{\dd{N}}{N} = \sigma (\theta) \dd{\Omega}
  \end{align}
  上式の左辺を言葉に直すと，
  \begin{align}
    \frac{\text{単位時間に位置$(r,\theta)$にある$\dd{S}$に入射する粒子数}}{\text{単位時間単位面積当たりの入射粒子数}}
  \end{align}
  である．これは
  \begin{align}
    \frac{\text{単位時間に位置$(r,\theta)$にある$\dd{S}$に粒子が入射する確率}}{\text{単位時間単位面積当たりに粒子が入射する確率}} = \frac{j_z \dd{S}}{j_z}
  \end{align}
  と等しいので
  \begin{align}
    \sigma(\theta)\dd{\Omega} = \frac{\abs{f(\theta)}^2\dd{S}}{r^2}
  \end{align}
  が得られる．よって，$\dd{\Omega} = \frac{\dd{S}}{r^2}$だから，
  \begin{align}
    \sigma(\theta) = \abs{f(\theta)}^2
  \end{align}
  という関係が成り立つ．これは，散乱振幅から微分断面積が求められることを意味する．
  \begin{itembox}[l]{散乱振幅と微分断面積の関係}
    \begin{equation}
      \sigma(\theta) = \abs{f(\theta)}^2
    \end{equation}
  \end{itembox}

  次に$f(\theta)$を求めるための式を作る．
  散乱のSchrödinger方程式は
  \begin{align}
    \qty(-\frac{\hbar^2}{2m}\nabla^2 + V(r))\psi(\bm{r}) = E\psi(\bm{r})
  \end{align}
  である．ここで，
  \begin{align}
    \begin{dcases}
      \kappa = \frac{\sqrt{2mE}}{\hbar}\\
      U(\bm{r}) = \frac{2m}{\hbar^2} V(\bm{r})
    \end{dcases}
  \end{align}
  とおくと，
  \begin{align}
    \label{Helmholtz equation}
    \qty(\nabla^2 + \kappa^2)\psi(\bm{r}) = U(\bm{r}) \psi(\bm{r})
  \end{align}
  と表せる．式(\ref{Helmholts equation})の解は，
  \begin{align}
    \qty(\nabla^2 + \kappa^2)\phi(\bm{r}) = 0
  \end{align}
  の一般解$\phi(\bm{r}) = \e^{\i k z}$と
  \begin{align}
    \label{Tokkai}
    \qty(\nabla^2 + \kappa^2)\chi(\bm{r}) = U(\bm{r}) \chi(\bm{r})
  \end{align}
  と特解$\chi(\bm{r})$の和である．
  
  では，式(\ref{Tokkai})の特解を求めよう．レシピはこうである．
  \begin{enumerate}
    \item $\qty(\nabla^2 + \kappa^2)G_0(\bm{r}) = \delta(\bm{r})$を満たすGreen関数$G_0$を求める．
    \item $\chi(\bm{r}) = \int \dd{\bm{r'}} G_0(\bm{r} - \bm{r'}) U(\bm{r'}) \phi(\bm{r'})$から特解を求める．
  \end{enumerate}
  2から特解が求められるのは，
  \begin{align}
    \qty(\nabla^2 + \kappa^2)\chi(\bm{r}) &= (\nabla^2 + \kappa^2)\int \dd{\bm{r'}} G_0(\bm{r} - \bm{r'}) U(\bm{r'}) \phi(\bm{r'})\\
    &= \int \dd{\bm{r'}} \delta(\bm{r} - \bm{r'}) U(\bm{r'}) \phi(\bm{r'})\\
    &= U(\bm{r}) \phi(\bm{r})
  \end{align}
  が成り立つためである．

  Green関数を求めよう．
  \begin{align}
    \qty(\nabla^2 + \kappa^2)G_0(\bm{r}) = \delta(\bm{r})
  \end{align}
  の両辺をFourier変換する．
  \begin{align}
    \qty(\nabla^2 + \kappa^2)\int \dd{\bm{k'}} G_0(\bm{k'}) \e^{\i\bm{k'}\cdot \bm{r}} &= \frac{1}{(2\pi)^3}\int \dd{\bm{k'}} \e^{\i\bm{k'}\cdot\bm{r}}\\
    \int \dd{\bm{k}} \qty(-\bm{k'}^2 + \kappa^2)G_0(\bm{k'}) \e^{\i \bm{k'}\cdot \bm{r}} &= \frac{1}{(2\pi)^3} \int \dd{\bm{k'}} \e^{\i \bm{k'}\cdot \bm{r}}
  \end{align}
  両辺を比較すると
  \begin{align}
    (\kappa^2 - k'^2)G_0(\bm{k'}) &= \frac{1}{(2\pi)^2}\\
    G_0(\bm{k'}) &= \frac{1}{(2\pi)^2 (\kappa^2 - k'^2)}
  \end{align}
  を得る．よってGreen関数は
  \begin{align}
    G_0(\bm{r}) = \frac{1}{(2\pi)^2}\int \dd{\bm{k'}} \frac{\e^{\i\bm{k'}\cdot\bm{r}}}{\kappa^2 - k'^2}
  \end{align}
  と表される．極座標に変換しこの積分を実行する．
  \begin{align}
    G_0(\bm{r}) &= \frac{1}{(2\pi)^2} \int \frac{\exp(\i k' r \cos \theta)}{\kappa^2 - k'^2}k'^2\sin\theta\dd{\theta}\dd{\phi}\dd{k'}\\
    &= \int \frac{1}{(2\pi)^3}2\pi \int_{0}^{\infty} k'^2 \dd{k'} \int_{0}^{\pi}\sin\theta\dd{\theta}\frac{\exp(\i k' r\cos\theta)}{\kappa^2 - k'^2}\\
    &= \frac{1}{(2\pi)^2} \int_{0}^{\infty} \frac{k'^2\dd{k'}}{\kappa^2 - k'^2}\frac{\e^{\i k'r} - \e^{-\i k'r}}{\i k'r}\\
    &= \frac{1}{2\pi^2\i r} \int_{0}^{\infty}\dd{k'} \frac{\e^{\i k'r} - \e^{-\i k'r}}{(\kappa - k')(\kappa + k')}\\
    &= -\frac{1}{8\pi^2\i r} \int_{-\infty}^{\infty} k' \dd{k'} \frac{\e^{\i k'r} - \e^{-\i k'r}}{(k' - \kappa)(k' + \kappa)}\\
    &= -\frac{1}{8\pi^2\i r} \kappa \qty{2\pi \i \frac{\e^{\i \kappa r} - \e^{-\i \kappa r}}{2 \kappa} + 2\pi \i \frac{\e^{- \i \kappa r} - \e^{\i \kappa r}}{-2 \kappa}}\\
    &= -\frac{1}{4\pi r}(\e^{\i \kappa r} - \e^{-\i \kappa r})
  \end{align}
  したがって，Green関数は外向きの球面波と内向きの球面波の重ね合わせになっている．物理的に意味のあるのは外向きの球面波であるため，こちらを採用する．

  以上の計算からSchrödinger方程式(\ref{SE_Scattering})の形式解は
  \begin{align}
    \psi(\bm{r}) = \e^{\i kz} - \frac{1}{4\pi} \int \frac{\e^{\i k \abs{\bm{r} - \bm{r'}}}}{\abs{\bm{r} - \bm{r'}}} \frac{2m}{\hbar^2}V(r')\psi(\bm{r'}) \dd{r'}
  \end{align}
  である．これは平面波と球面波の和となっている．
\end{document}
